\documentclass[11pt]{article}
\usepackage[margin=0.85in]{geometry}
\usepackage{amsmath,booktabs,xurl,hyperref}
\usepackage[T1]{fontenc}
\setlength{\emergencystretch}{2em}
\setlength{\tabcolsep}{4pt}
\title{Root Causes of the Real-Time-Iteration Controller's Scenario Failures}
\author{Pan Dynamics Collision Avoidance Research}
\date{October 1, 2026}
\begin{document}
\maketitle

\section{Summary}
The fourteen collision-threat scenarios were rerun on the real-time-iteration
(RTI) controller of commit \path{5f7497718d4774a2515330650ee8815ccfa67b20}.
Three cases recover, two collide, and nine end the 80-s window far from the
given path. No control call was aborted. In all 20{,}684 frames the
safety (PCBF) stage reported an optimum of at most $3.9\times10^{-6}$, so the
optimizer never stated that a safety relaxation was needed. The failures
therefore do not come from the solver being unable to return a result. They
come from what the returned result is allowed to mean:
\begin{enumerate}
\item \textbf{Collisions: separation is enforced only at sample times.}
Collision rows exist at the start and the interpolated midpoint of each 50-ms
hold, with a 6-mm margin. In both collisions every sampled clearance is at
least 9.7~mm and the one-step model error is below 10~$\mu$m, yet the
continuous motion overlaps the target between samples for 15.4 and 2.7~ms
(depth 29.4 and 4.8~mm) at about 20~m/s relative speed.
\item \textbf{Non-recovery: the RTI trust region is centered on a plan that the
objectives do not determine.} The second stage minimizes a one-step CLF
slack, which depends only on the first input. The later inputs are set by
the interior-point solution inside the trust region around the previous plan,
and they cancel the first-step correction. Shifting that plan re-centers the
next frame's $\pm0.075$-rad steering box near the old input, so the issued
steering stays at a box edge in 82--99\% of late frames and the vehicle cannot
turn back. The same mechanism in the braking channel produces a speed limit
cycle in the two 15-m/s head-on cases. Far from the path, no admissible input can satisfy the one-step
CLF decrease, so the CLF stage cannot break the cycle by itself.
\item \textbf{Warm-start infeasibility: affine plans re-rolled through the
nonlinear model violate hard tail or terminal rows.} This happened three
times and was repaired each time by the flow restart.
\end{enumerate}
Diagnostic counterfactuals confirm each mechanism. Six extra collision samples
per hold, or a 5-cm margin, remove both contacts. A third stage that gives the
whole plan a tracking objective recovers nine of ten unrecovered runs, but five
of the nine then touch the target between samples or through relaxation. The
counterfactuals are not validated controller changes, and no production file was
modified.

\section{Campaign}
Each case uses \texttt{runNonlinearPredictiveSafetyValidation} with the
original fixture, 8 or 15~m/s with prefix 8 or 16, \texttt{Frames=1600},
\texttt{RequireCollisionThreat=true} and a five-second recovery dwell. Ego
measurements are exact; the target motion is known; there are no road
constraints, observer, noise or random draws. The plant is ODE45 (relative and
absolute tolerances $10^{-11}$ and $10^{-12}$), and clearance is sampled at 31
nodes per hold. The campaign ran on a frozen copy whose controller,
configuration, driver and native kernels are byte-identical to
\path{5f74977}. Its issued inputs agree bitwise, frame by frame, with the
independent run reported in \path{RTI_FLOW_REINITIALIZATION_20261001.tex}.
\begin{center}\footnotesize
\begin{tabular}{rlrrrrrl}
\toprule
Speed & Scenario & Time (s) & Min gap (m) & Final $e_y$ (m) & CLF $-7$ & Restarts & Outcome\\
\midrule
8 & headOn & 49.00 & 0.583 & 0.0 & 130 & 0 & recovered\\
8 & acceleratingHeadOn & 52.30 & 0.368 & 0.0 & 122 & 0 & recovered\\
8 & brakingLead & 52.90 & 1.060 & 0.0 & 111 & 0 & recovered\\
8 & crossing & 80.00 & 0.699 & $-621.1$ & 40 & 0 & not recovered\\
8 & turningCrossing & 80.00 & 0.411 & $-124.4$ & 13 & 0 & not recovered\\
8 & curvedHeadOn & 80.00 & 0.664 & $-72.6$ & 143 & 0 & not recovered\\
8 & curvedCrossing & 80.00 & 0.545 & 25.5 & 12 & 0 & not recovered\\
15 & headOn & 80.00 & 0.000 & 0.0 & 30 & 0 & collision; speed cycle\\
15 & acceleratingHeadOn & 80.00 & 0.000 & 0.0 & 41 & 0 & collision; speed cycle\\
15 & brakingLead & 80.00 & 0.655 & $-232.5$ & 95 & 0 & not recovered\\
15 & crossing & 80.00 & 0.381 & $-627.4$ & 166 & 0 & not recovered\\
15 & turningCrossing & 80.00 & 0.081 & 471.6 & 131 & 0 & not recovered\\
15 & curvedHeadOn & 80.00 & 0.780 & $-26.2$ & 307 & 1 & not recovered\\
15 & curvedCrossing & 80.00 & 0.153 & $-139.9$ & 152 & 2 & not recovered\\
\bottomrule
\end{tabular}\end{center}
``CLF $-7$'' counts frames whose second stage returned coneprog exit $-7$; the
finite point is still used, as the controller's output contract specifies.
The PCBF stage never returned $-7$. Across the campaign, 1{,}493 of 20{,}684
frames (7.2\%) have this exit. The two collision cases end on the given path
but never satisfy the five-second dwell because of a speed limit cycle
(Section~\ref{sec:speed}).

\section{Collisions occur between constrained samples}
Both collisions were reproduced by an instrumented replay whose issued inputs
match the campaign exactly. The table measures the colliding hold with a
5{,}001-point ODE45 solution of the issued input.
\begin{center}\footnotesize
\begin{tabular}{lrr}
\toprule
 & headOn & acceleratingHeadOn\\
\midrule
Hold start (s) & 1.10 & 1.05\\
PCBF optimum at this call / maximum up to it & $4.2\times10^{-8}$ / $3.2\times10^{-7}$ & $9.8\times10^{-8}$ / $3.6\times10^{-7}$\\
Clearance at start / true midpoint / end (mm) & 9.72 / 34.21 / 166.41 & 49.05 / 29.41 / 168.40\\
Clearance at the controller's interpolated midpoint (mm) & 34.53 & 29.15\\
One-step RK4 position error (m) & $4.0\times10^{-6}$ & $9.7\times10^{-7}$\\
Strict overlap (ms after hold start) & 3.12--18.51 & 17.05--19.71\\
Maximum penetration depth (mm) & 29.35 & 4.76\\
Offline interval certificate & rejected & rejected\\
Distance Lipschitz bound (m/s) & 20.69 & 23.20\\
\bottomrule
\end{tabular}\end{center}
Every sample the controller constrains clears the 6-mm margin, the solver
reports two positive exits, and the prediction of the issued hold is accurate
to micrometres. The vehicles' corners pass each other in a few milliseconds
while the distance is sampled every 25~ms. With the measured Lipschitz bounds,
clearance can fall by up to 0.26--0.29~m below the sampled values within
that spacing, which is far more than the 6-mm margin. The existing offline
function \texttt{predictiveSafetyGeometry.intervalClearance}, applied to each
hold's node-to-node pose interpolant, rejects both holds. Its online use was
removed in commit \path{905c8bd}; the current formulation has no
inter-sample constraint.

Counterfactual runs over the first 4~s of both encounters, with everything
else unchanged:
\begin{center}\footnotesize
\begin{tabular}{lrr}
\toprule
Variant & headOn min gap (m) & acceleratingHeadOn min gap (m)\\
\midrule
Original (6-mm margin, 2 samples per hold) & 0 (overlap) & 0 (overlap)\\
Six extra interpolated collision samples per hold & 0.114 & 0.054\\
Margin 0.05~m & 0.035 & 0.031\\
Margin 0.10~m & 0.095 & 0.086\\
Margin 0.30~m & 0.286 & aborted at 0.15~s (no PCBF result)\\
\bottomrule
\end{tabular}\end{center}
Denser sampling or a few centimetres of margin remove these two contacts. The
largest margin also makes one encounter unsolvable from its warm start. These
runs identify the mechanism; they do not certify a margin or sampling density.

\section{Non-recovery: the trust region is centered on an undetermined plan}
\subsection{What the issued inputs look like}
Instrumented replays store each frame's linearization anchor, which is the
previous plan's inputs shifted by one hold and rolled out from the measured
state. The input correction box is $|u_i-\bar u_i|\le(0.075,\,0.125)$ around
that anchor. After 10~s (1{,}400 frames per case):
\begin{center}\footnotesize
\begin{tabular}{rlrrrrr}
\toprule
Speed & Scenario & Steering at box edge & Next input counter-corrected & In encounter & Median $e_v$ (m/s)\\
\midrule
8 & crossing & 1367 & 1400 & 1400 & 0.00\\
8 & turningCrossing & 1335 & 1342 & 122 & $-1.22$\\
8 & curvedHeadOn & 1252 & 1400 & 0 & $-1.08$\\
8 & curvedCrossing & 1151 & 1285 & 396 & 0.42\\
15 & brakingLead & 1333 & 1228 & 349 & $-3.00$\\
15 & crossing & 1379 & 1392 & 1400 & $-7.01$\\
15 & turningCrossing & 1324 & 1340 & 85 & $-3.01$\\
15 & curvedHeadOn & 1336 & 1373 & 69 & $-1.50$\\
15 & curvedCrossing & 1289 & 1340 & 310 & $-2.64$\\
\bottomrule
\end{tabular}\end{center}
In the nine unrecovered cases the issued steering lies on a box edge in
82--99\% of these frames. In 97--100\% of those edge frames it is the edge that
reduces the CLF's heading term, that is, the side that would turn the vehicle
back. The second planned input is corrected in the opposite direction in
88--100\% of frames. In 8-m/s crossing, for example, the issued steering is
within 1~mrad of the top of its box in all 1{,}520 frames after 4~s, while the
anchor steering stays between $-0.144$ and $-0.056$~rad. The issued steering
therefore never exceeds 0.019~rad, and the heading error stays between
$-1.58$ and $-1.54$~rad from 8~s to 80~s.

\subsection{Why the anchor never moves}
Single calls were rebuilt from captured controller inputs with a copy of the
controller's local functions. All rebuilt calls reproduce the issued input
exactly. Modified copies of each call isolate the mechanism:
\begin{center}\footnotesize
\begin{tabular}{lrrrrrr}
\toprule
 & \multicolumn{2}{c}{CLF slack} & \multicolumn{2}{c}{$\sum$ later corr.} & \multicolumn{2}{c}{Free $\delta_0$}\\
Case, time & original & future fixed & original & no state box & $\delta_0$ & actual $\Delta V$\\
\midrule
8 crossing, 4~s & 1140 & 1144 & $-0.080$ & $+0.685$ & 0.435 & $+827$\\
8 crossing, 8~s & 4829 & 4840 & $-0.085$ & $+0.750$ & 0.411 & $+2797$\\
15 crossing, 20~s & 21160 & 21180 & $-0.076$ & $+0.592$ & 0.379 & $+6562$\\
8 turningCrossing, 30~s & 13320 & 13340 & $-0.059$ & $-0.010$ & 0.451 & $+3964$\\
8 curvedHeadOn, 30~s & 606.5 & 606.6 & $-0.062$ & $-0.013$ & 0.443 & $+200$\\
\bottomrule
\end{tabular}\end{center}
Steering values are in radians; ``later corr.'' is the sum of the steering
corrections of all inputs after the first. ``Future fixed'' holds every input
after the first at its anchor value. ``No state box'' removes the state
correction bounds. ``Free $\delta_0$'' removes only the steering bound on the
issued input and reports the steering chosen.
Four facts follow:
\begin{itemize}
\item \emph{The objectives do not determine the future plan.} Fixing every
later input at its anchor changes the minimized CLF slack by at most 0.35\%.
The returned later inputs are one point of a large optimal set.
\item \emph{The returned future cancels the first-step correction.} The first
steering moves $+0.075$~rad, and the later corrections sum to $-0.059$ to
$-0.085$~rad. The planned endpoint heading stays within 4~mrad of the anchor's.
The first counter-steer is $-0.0185$~rad at 8~s in 8-m/s crossing, and the
counter-steer decays over about 13 stages.
\item \emph{The state trust box is the main contributor in these calls.}
Without it, the crossing cases' later corrections become $+0.59$ to $+0.75$~rad
and the plans turn toward the path; in the other two calls the cancellation
shrinks to $-0.010$ and $-0.013$~rad. In the original problem the state box is
only 15\% used at the solution, so the effect is interior-point centering, not
an active constraint. In a closed loop without state boxes a smaller
cancellation remains ($-0.081$~rad at 30~s in 8-m/s crossing).
\item \emph{The issued input wants to turn back.} With only $\delta_0$ freed,
the optimizer chooses 0.38--0.45~rad toward the path. The later plan again
cancels it, by $-0.39$ to $-0.48$~rad.
\end{itemize}
The next frame's anchor is this plan shifted by one hold. Its first input is
the counter-corrected second input, so the new box again excludes the turning
action. The loop is self-sustaining.

\subsection{The CLF cannot be decreased far from the path}
The CLF is the LQR quadratic $V=e^TPe$ of the nominal trim. At 8~m/s,
$P_{11}=76.0$, $P_{12}=148.7$ and $P_{22}=1564$ in lateral metres and heading
radians. At 8~s in 8-m/s crossing ($e_y=-44.8$~m, $e_\psi=-1.54$~rad), the
vehicle moves 0.40~m further from the path in each hold whatever it steers.
This alone raises $V$ by $2(P_{11}e_y+P_{12}e_\psi)\,\Delta e_y\approx+2906$. The
required 1\% decrease is $-1768$. Freeing the first input gives predicted
changes of $+2398$ to $+2474$ and actual changes of $+2797$ to $+3059$. Only removing every trust
bound gives a predicted decrease, and it does so with steering of 2.1--3.4~rad,
outside the Fiala slip domain; ODE45 cannot integrate it. The CLF's preferred
heading at this lateral error, $-(P_{12}/P_{22})e_y=4.3$~rad, is not a valid
heading. The second stage can supply only a greedy one-step direction. It cannot
certify or plan the multi-second turnaround these states need, so the slack
stays positive and the trust-region loop decides the outcome.

\subsection{Two scenario-specific consequences}
\paragraph{Escort capture in the crossing cases.} Both crossing avoidances turn
the ego into the target's direction of travel. Because the ego cannot turn back,
from 8~s to 80~s it drives beside the target at the target's speed and heading.
The center distance is 7.2--7.4~m at 8~m/s and 5.2--6.2~m at 15~m/s. At
15~m/s the ego is held near the target's 8~m/s from 16~s on, a speed error of
$-7$~m/s. The 30-m encounter never ends (all 1{,}400 late
frames), so every collision and terminal-encounter row stays active for the
rest of the run.

\paragraph{Speed limit cycle.}\label{sec:speed} In the two 15-m/s head-on cases
the lateral errors vanish, but the speed error cycles between 0 and about
$-0.37$~m/s with a period of about 10~s. At 68.35~s in head-on, the braking
anchor has drifted to $-0.124$ while cruise needs about $+0.02$ to $+0.04$. The
issued braking ratio is pinned at the box edge, $0.0006$, so the vehicle
coasts. Removing only the first-input box gives braking $0.039$ and an actual
CLF decrease. Fixing the future at the drifted anchor makes the result worse
(CLF slack 0.083, exit $-7$). The braking channel has the same trap.

\subsection{Causal tests in closed loop}
Diagnostic copies of \texttt{solvePredictiveControl.m} were placed ahead of the
frozen controller on the MATLAB path:
\begin{center}\footnotesize
\begin{tabular}{p{0.42\linewidth}p{0.5\linewidth}}
\toprule
Variant & Result\\
\midrule
No steering bound on the issued input, activated at 0~s, at 4~s, or when the
encounter ends &
Solved steering leaves the Fiala slip domain within 0--0.9~s of activation
(1.58~rad in the plant at 0.50~s in 8-m/s crossing); 7 of 7 activated runs abort\\
First-input box centered on the previously applied input, same activations &
No PCBF result within 0.5--8~s of activation; 7 of 7 activated runs abort\\
No state trust box, from 0~s &
Both crossing cases still escort the target to 80~s; 8-m/s turningCrossing
recovers laterally but ends with a $+0.99$~m/s speed error\\
Third stage minimizing horizon tracking cost subject to both achieved slack levels &
9 of 10 runs recover by 18.05~s; 5 of those 9 touch the target (table below)\\
\bottomrule
\end{tabular}\end{center}
Removing the box fails because the affine model is then used far outside its
validity; this is consistent with \path{UNBOUNDED_STEERING_20261001.tex}.
Moving only the first input's box makes it inconsistent with the rest of the
plan's anchor. Neither result argues for the current center: the box is needed,
but around a plan that heads back.

The third stage gives every planned input an objective without changing either
achieved slack level. It was run on the nine laterally unrecovered cases and on
15-m/s head-on:
\begin{center}\footnotesize
\begin{tabular}{rlrrl}
\toprule
Speed & Scenario & Recovered at (s) & Min gap (m) & Contact\\
\midrule
8 & crossing & 13.00 & 0 & between samples, 1.85-s hold\\
8 & turningCrossing & 15.80 & 0 & positive PCBF relaxation 0.0106, 1.80-s hold\\
8 & curvedHeadOn & 13.00 & 0.031 & none\\
8 & curvedCrossing & 18.05 & 0 & between samples, 1.80- and 12.45-s holds\\
15 & headOn & 13.00 & 0.0056 & none\\
15 & brakingLead & --- & 0.753 & aborted at 4.00~s, no PCBF result\\
15 & crossing & 13.00 & 0.079 & none\\
15 & turningCrossing & 13.25 & 0 & between samples, 1.70-s hold\\
15 & curvedHeadOn & 13.00 & 0.0061 & none\\
15 & curvedCrossing & 13.00 & 0 & between samples, 1.80-s hold\\
\bottomrule
\end{tabular}\end{center}
Thirteen seconds is the earliest possible confirmation. Nine of the ten runs
recover; the original controller recovers none of them, and the 15-m/s head-on
speed cycle disappears. This confirms that the undetermined future plan is the
cause of the non-recovery. It also shows why a recovery objective alone is not
a fix. Five of the nine recovering runs touch the target, because the plans now
try to return to the path while the target is still present. Four contacts occur between samples with PCBF optima below
$1.1\times10^{-7}$, the mechanism of the two campaign collisions. One plan accepts
a positive PCBF relaxation. The 8-m/s curvedHeadOn run later meets the same target
on the curve again, 6.6~m away at 80~s, and deviates to avoid it. This
diagnostic stage reached the 5-s call budget in 0--25 calls per run and used those
finite points, so its details can depend on machine load.

\section{Warm-start infeasibility}
The three flow restarts occurred when the warm PCBF problem returned exit $-2$.
Rebuilt calls identify the conflicting rows:
\begin{center}\footnotesize
\begin{tabular}{lll}
\toprule
Call & Feasible after removing & Geometry of the shifted rollout\\
\midrule
15 curvedHeadOn, 0.10~s & hard tail collision rows, or all trust boxes & overlaps the target at tail stage 26\\
15 curvedCrossing, 19.80~s & terminal endpoint cone, or all trust boxes & endpoint far outside the core; 4.5-m drift\\
15 curvedCrossing, 21.60~s & endpoint cone, state boxes, or all trust boxes & endpoint outside the core; 0.8-m drift\\
\bottomrule
\end{tabular}\end{center}
Removing only input boxes, only tail rows, only terminal clearance rows or only
midpoint state rows does not restore feasibility, except as listed. At 0.00 and 0.05~s
the target is 31.8 and 30.6~m away, beyond the 30-m range, so those problems
contain no collision rows; both returned plans pass through the target at the
same tail time (exact rectangle distance zero). At 0.10~s the target is 29.5~m
away, the shifted plan meets hard tail rows that it violates, and the trust
boxes cannot remove the overlap in one step. At 19.80~s the previous plan's endpoint
was inside the core in the affine model, but its nonlinear re-rollout drifts
4.5~m and ends with core membership 1196 against a squared radius of
$9.5\times10^{-7}$. Hard tail and terminal rows combined with fixed trust boxes
cannot absorb this one-step mismatch. The flow restart repaired all three
calls.

\section{Root causes and implications}
\begin{enumerate}
\item \emph{Safety is sampled, not continuous.} Two collision samples per hold
with a 6-mm margin cannot bound inter-sample contact at 20~m/s relative speed.
An inter-sample certificate, or a margin tied to the distance Lipschitz bound
and the sample spacing, is needed. More sampling alone narrows but does not
close the gap.
\item \emph{The future plan has no objective, but it defines the next trust
region.} The two-stage objectives determine the prefix slacks and the first
input only. The interior-point solution returns later inputs that undo the
first-step correction, and shifting them pins the next issued input. Any fix
must give the whole plan a recovery-consistent objective, or center the local
model on such a plan. Deleting or enlarging the boxes alone leaves the model
invalid.
\item \emph{The local one-step CLF does not reach.} Far from the path the 1\%
decrease is unattainable for any admissible input, so the CLF cannot drive the
turnaround without a multi-step mechanism.
\item \emph{Long affine plans are brittle at the hard tail and tiny terminal
core.} One-step model mismatch accumulates over 3.8~s, and the core radius is
$9.8\times10^{-4}$. This produces the warm infeasibilities that the flow
restart currently absorbs.
\end{enumerate}
The diagnostic third stage also shows that fixing the second cause makes the
first more visible: plans that return during an encounter press against the
sampled collision rows. The causes should be addressed together.

\section{Limits}
All results are nominal-model closed loops with exact state and a known
target. The 30-m range, the recovery criterion and the 80-s window are declared
conventions. Single-call ablations change one condition at a time from
reproduced controller inputs. They are not alternative closed-loop
trajectories. Counterfactual controllers are diagnostic copies run outside the
repository; none is proposed as a validated algorithm, and no production file
was changed. Several analyses ran concurrently with unrelated jobs, so wall
times are not benchmarks; the controller's 5-s budget was not reached in any
reported call. ``Not recovered'' means not within 80~s, not proved divergent.

\section{Artifacts}
\path{report/RTI_FAILURE_ROOT_CAUSES_20261001/} contains the outcome, collision,
trust-box, ablation and counterfactual tables; the diagnostic methods; patches
for every diagnostic controller copy; reproduction notes; and a manifest of
source and raw-artifact hashes. Full campaign JSON/MAT, replays, captured calls
and counterfactual records remain in
\path{/home/zai/.cache/collisionAvoidance/rti-failure-analysis-20261001/}.
\end{document}
